% Generated draft: requires mathematical and source review.
\documentclass[11pt]{article}
\usepackage[margin=25mm]{geometry}
\usepackage{fontspec}
\setmainfont{Latin Modern Roman}
\usepackage{amsmath,amssymb,mathtools,unicode-math}
\setmathfont{Latin Modern Math}
\usepackage{enumitem,needspace,xurl,hyperref}
\setlength{\parindent}{0pt}
\setlength{\parskip}{5pt}
\setlength{\emergencystretch}{3em}
\newcommand{\N}{\mathbb{N}}
\newcommand{\Z}{\mathbb{Z}}
\newcommand{\Q}{\mathbb{Q}}
\newcommand{\R}{\mathbb{R}}
\newcommand{\C}{\mathbb{C}}
\newcommand{\F}{\mathbb{F}}
\newcommand{\A}{\mathbb{A}}
\newcommand{\PP}{\mathbb{P}}
\newcommand{\E}{\mathbb{E}}
\newcommand{\eps}{\varepsilon}
\newcommand{\dd}{\,\mathrm{d}}
\begin{document}
\section*{UnsolvedMath ID: 30006428}
\textbf{Source problem code:} catalogue entry
\subsection{Definitions and mathematical statement}
A complex \(C^*\)-algebra \(A\) is simple if it is nonzero and its only closed two-sided ideals are \(\{0\}\) and \(A\); it is separable if it has a countable norm-dense subset. For \(a\in A\), let \(C^*(a)\) denote the smallest norm-closed \(*\)-subalgebra of \(A\) containing \(a\). Say that \(A\) is singly generated if \(A=C^*(a)\) for some \(a\in A\); the generator need not be self-adjoint.

The first question is whether
\[
\forall A\;\bigl[A\text{ is a unital simple separable complex }C^*\text{-algebra}\ \Longrightarrow\ \exists a\in A\; A=C^*(a)\bigr].
\]
The supplied dated curation reports a counterexample to this assertion; that report is not independently verified here.

Let \(F_2=\langle s,t\mid\ \rangle\) be the free group on two generators. For a discrete group \(\Gamma\), its left regular representation on \(\ell^2(\Gamma)\) is defined by \(\lambda_\Gamma(g)\delta_h=\delta_{gh}\), and its reduced group \(C^*\)-algebra is
\[
C_r^*(\Gamma)=\overline{\operatorname{span}\{\lambda_\Gamma(g):g\in\Gamma\}}^{\|\cdot\|}\subseteq B(\ell^2(\Gamma)).
\]
Using the standard meaning of a hyperbolic group, namely a finitely generated group whose Cayley graph for a finite generating set is Gromov hyperbolic, the group-determination question is whether
\[
\forall\Gamma\;\bigl[\Gamma\text{ is a hyperbolic group and }C_r^*(\Gamma)\cong C_r^*(F_2)\text{ as }C^*\text{-algebras}\ \Longrightarrow\ \Gamma\cong F_2\text{ as groups}\bigr].
\]
No torsion-free hypothesis is supplied. The source also asks about a possibly stronger \(C_r^*\)-superrigidity property, but does not define that property or specify restrictions on the algebra isomorphism. Consequently, the displayed implication formalizes only the group-determination question; the additional superrigidity target remains unresolved at the level of formulation.
\subsection{Short English statement}
The questions ask whether every unital simple separable \(C^*\)-algebra has one generator, and whether a hyperbolic group with reduced group \(C^*\)-algebra isomorphic to that of \(F_2\) must itself be isomorphic to \(F_2\). The supplied curation reports a negative answer to the first question and describes the second as open. The stronger intended meaning of superrigidity is unspecified.
\subsection{Sources}
\begin{itemize}
\item Local UnsolvedMath registry, permanent ID \(30006428\), source code \texttt{catalogue entry}, title: Single Generation and Free-Group \(C^*\)-Superrigidity. The supplied record lists proposed year 2025 and identifies no individual proposer.
\item Supplied report citation: \emph{C*-Algebras (2026)}, Oberwolfach Reports, DOI \texttt{10.4171/owr/2025/35}, pp.\ 1887--1948. \url{https://doi.org/10.4171/owr/2025/35}. This is the supplied workshop-report provenance; the year discrepancy between the title, DOI, and proposed year is unresolved.
\item Supplied literature source, labeled Elliott--Li--Niu counterexample: \url{https://arxiv.org/abs/2608.13707}. The supplied assessment attributes to an August 2026 construction a simple unital separable AH (approximately homogeneous) algebra that is not singly generated.
\item Supplied literature source, labeled Oberwolfach Report 35: \url{https://doi.org/10.4171/owr/2025/35}.
\item Literature assessment date: \texttt{2026-08-22}. This is supplied dated curation, not independently verified. It labels the combined record partially solved, reports the universal single-generation assertion disproved, and describes the separate \(C^*\)-superrigidity question for \(F_2\) within hyperbolic groups as open. The supplied background explicitly requests independent verification of these time-sensitive statuses and cited claims.
\end{itemize}
\end{document}
